\chapter{Experimental Setup}\label{ch:setup}

\section{Tasks and datasets}

Three datasets were recorded on the SO-101 by teleoperation at 30~fps (Table~\ref{tab:datasets}). The idle start and
end of every episode were trimmed automatically. In each dataset a fixed, randomly chosen (seeded) set of episodes is
held out: it is never used for training and serves the offline evaluation (Section~\ref{sec:offline}). The first
cube-in-box dataset was an initial, smaller recording on an earlier camera arrangement; the second repeats the task
with more demonstrations. The sorting task is the main task of this report: a coloured cube must be placed in the
square of the same colour.

\begin{table}[htbp]
  \centering
  \caption{Datasets. Joint units: degrees, or LeRobot's normalised range for the earlier recordings.}
  \label{tab:datasets}
  \small
  \gentable{tab_datasets}
\end{table}

\begin{figure}[htbp]
  \centering
  \begin{subfigure}[t]{0.48\textwidth}
    \centering
    \photo[\textwidth]{task_sort_one.jpg}{The sorting task: the cube and the coloured squares, seen by the scene camera.}
    \caption{Sorting.}
  \end{subfigure}\hfill
  \begin{subfigure}[t]{0.48\textwidth}
    \centering
    \photo[\textwidth]{task_cubebox.jpg}{The cube-in-box task.}
    \caption{Cube in box.}
  \end{subfigure}
  \caption{The two tasks.}
  \label{fig:tasks}
\end{figure}

\section{Policies and training}\label{sec:policies}

Seven policies were chosen to span small task-specific models and large VLAs: \smolvla, \groot, GR00T~N1.7,
\pizerofive, MolmoAct~2, ACT and Diffusion Policy. Each is trained with its authors' published recipe, unchanged
(Appendix~\ref{app:recipes}), so that differences between them come from the models and not from our tuning. Training
ran on the H200 and A100 servers; every run was logged, and every checkpoint was stored. The checkpoint served to the
robot is the one with the lowest loss on the held-out episodes.

\begin{table}[htbp]
  \centering
  \caption{Training status on 20 November 2026 (trained: on the model hub; training: running; planned: not started).}
  \label{tab:training}
  \small
  \begin{tabular}{l c c c}
    \toprule
    Policy & Cube in box v0 & Cube in box v1 & Sorting \\
    \midrule
    \smolvla & trained & trained & trained \\
    \groot & trained & trained & trained \\
    ACT & -- & training & trained \\
    Diffusion Policy & -- & training & trained \\
    GR00T N1.7 & -- & planned & planned \\
    \pizerofive & -- & planned & planned \\
    MolmoAct 2 & -- & planned & training \\
    \bottomrule
  \end{tabular}
\end{table}

\section{Conditions}\label{sec:conditions}

Four conditions are compared (Table~\ref{tab:conditions}). The network is the real WiFi link in all of them; conditions
2--4 add the fault mix of Figure~\ref{fig:faultmix}. Condition~2 against~1 measures the cost of the faults, 3
against~2 whether \drtc{} recovers it, and 4 against~3 what smoothing adds.

\begin{table}[htbp]
  \centering
  \caption{The four conditions.}
  \label{tab:conditions}
  \begin{tabularx}{\textwidth}{c l l X}
    \toprule
    \# & Condition & Network & Robot side \\
    \midrule
    1 & Baseline & normal & LeRobot remote inference as shipped: no \drtc, no filter, ends the run after 1~s
      without actions \\
    2 & Baseline + faults & fault mix & as 1 \\
    3 & \drtc{} + faults & fault mix & \drtc{} with our fixes; no filter; holds the arm through stalls \\
    4 & \drtc{} + smoothing + faults & fault mix & as 3, with the Butterworth filter and adaptive playback \\
    \bottomrule
  \end{tabularx}
\end{table}

\begin{table}[htbp]
  \centering
  \caption{Which policies are evaluated under which conditions.}
  \label{tab:plan}
  \small
  \begin{tabularx}{\textwidth}{l X l}
    \toprule
    Dataset & Policies & Conditions \\
    \midrule
    Sorting & \smolvla, \groot & 1--4 (the main experiment) \\
    Sorting & all seven & 1 \\
    Cube in box v1 & all seven & 1 \\
    Cube in box v0 & \smolvla, \groot & 1 \\
    Bench (replay robot) & \smolvla, \groot{} on sorting & 1--4, with each fault type alone and the mix \\
    \bottomrule
  \end{tabularx}
\end{table}

\section{Protocol}

\begin{itemize}
\item \textbf{Trials.} Ten valid trials per policy, condition and task, in rounds with a seeded random order
      (Section~\ref{sec:evalsys}). The operator chooses the object layout of each round; all trials of a round use it.
\item \textbf{Success.} The operator presses the success key before the time limit (30~s for sorting) when the task's
      written success rule is met. For sorting: \emph{the cube is fully inside the square of its colour and the gripper
      has let go of it.}
\item \textbf{Faults.} A trial ended by a terminal fault after the arm has started moving is a failure. A fault before
      the first action is a setup problem: that trial is repeated.
\item \textbf{Validity.} A trial is voided only for a setup mistake (an object not on its mark, a cable caught), never
      for a policy's failure; voided trials are kept and listed.
\item \textbf{Machines.} For every robot experiment the laptop was the robot computer and the A100 the model server.
\end{itemize}

\section{Metrics}\label{sec:metrics}

Table~\ref{tab:metrics} summarises the metrics; Appendix~\ref{app:metrics} gives the full definitions. Every metric is
first reduced to one value per trial and then summarised over the trials (median and interquartile range), so long
trials do not weigh more than short ones. Times are measured on one machine's monotonic clock and never across
machines.

\begin{table}[htbp]
  \centering
  \caption{Metrics.}
  \label{tab:metrics}
  \small
  \begin{tabularx}{\textwidth}{l l X}
    \toprule
    & Metric & Definition \\
    \midrule
    M1 & Success rate & Successful trials over valid trials, with a Wilson 95\% interval \\
    M1f & Ended by a fault & Trials that a terminal fault ended after the first action \\
    M1t & Time to success & From the first action to the success key \\
    M2 & Stall fraction & Share of control ticks after the first action without a new action \\
    M3 & Stalls & Number of stall events; the longest stall in seconds \\
    M4 & Turnaround & From sending an observation to receiving its chunk (median and 95th percentile) \\
    M5 & Estimate error & The scheduler's latency estimate minus the measured turnaround \\
    M6 & Inference time & The server's reported time per chunk; $\mathrm{M4}-\mathrm{M6}$ is the network share \\
    M7 & Jerk & Root mean square of the third time derivative of the commanded arm joints \\
    M8 & Chunk switch, largest step & Disagreement between chunks at a switch; the largest change between two ticks \\
    M9 & Held-out error & Offline: mean absolute error of the first 15 predicted actions on held-out episodes \\
    M10 & Inference cost & Offline: inference time over 200 requests, and peak GPU memory \\
    \bottomrule
  \end{tabularx}
\end{table}

With $T$ the ticks after the first action and $s_t=1$ when tick $t$ has no new action, the stall fraction is
$\mathrm{M2} = \frac{1}{T}\sum_t s_t$. With $q_{j,t}$ the commanded position of arm joint $j$ on the 30~Hz tick grid
(held through stalls), the jerk is
\begin{equation}
  \mathrm{M7} = \sqrt{\frac{1}{5\,(T-3)} \sum_{j=1}^{5} \sum_{t} \Bigl( f^3\,\bigl(q_{j,t+3} - 3q_{j,t+2} + 3q_{j,t+1}
  - q_{j,t}\bigr) \Bigr)^{2}}, \qquad f = 30~\mathrm{Hz},
\end{equation}
reported in \jerkunit. Lower is smoother.

\section{Statistical analysis}

Success rates are reported with Wilson 95\% intervals~\cite{wilson1927}. Differences in success between two
conditions of the same policy are tested with Fisher's exact test~\cite{fisher1935}, and differences in the other
metrics with the two-sided Mann--Whitney $U$ test~\cite{mann1947} on the per-trial values. Comparisons are made only
within one evaluation (same sessions, interleaved trials). We report $p$-values without correction for multiple
comparisons and read them with the effect sizes; Section~\ref{sec:limits} discusses this.
